\documentclass[11pt,a4paper]{article}

\usepackage[utf8]{inputenc}
\usepackage[T1]{fontenc}
\usepackage[ngerman]{babel}
\usepackage{amsmath,amssymb}
\usepackage{geometry}
\usepackage{hyperref}
\usepackage{setspace}
\usepackage{xcolor}
\usepackage{titlesec}
\usepackage{enumitem}
\usepackage{microtype}
\usepackage{fancyhdr}

% Layout
\geometry{margin=2.5cm,headheight=15pt}
\onehalfspacing
\setlength{\parindent}{0pt}
\setlength{\parskip}{0.8em}

% Farben
\definecolor{sectioncolor}{RGB}{0,51,102}
\definecolor{textcolor}{RGB}{30,30,30}

% Titelformatierung
\titleformat{\section}
  {\Large\bfseries\color{sectioncolor}}
  {}
  {0em}
  {}
  [\vspace{-0.3em}\rule{\textwidth}{0.8pt}\vspace{0.5em}]

\titleformat{\subsection}
  {\large\bfseries\color{sectioncolor}}
  {}
  {0em}
  {}

% Kopfzeile
\pagestyle{fancy}
\fancyhf{}
\fancyhead[L]{\small\textit{Von Platon zur Peanokugel}}
\fancyhead[R]{\small\thepage}
\renewcommand{\headrulewidth}{0.4pt}

% Hyperref-Einstellungen
\hypersetup{
    colorlinks=true,
    linkcolor=sectioncolor,
    urlcolor=sectioncolor,
    citecolor=sectioncolor
}

\title{\textbf{Von Platon zur Peanokugel:}\\
\large Wie man ein geometrisches Gedankenexperiment\\
in eine numerische Physik-Testbank gießt}
\author{}
\date{\today}

\begin{document}

\maketitle

\vspace{1em}

\begin{center}
\small
\textit{(eine Spektrum-der-Wissenschaft-Anmutung über Ziele, Ergebnisse\\
und Voraussetzungen unserer heutigen BM-Session)}
\end{center}

\vspace{2em}

\noindent Es beginnt wie bei Einstein: nicht mit einem Messgerät, sondern mit einem Bild im Kopf. Zwei Körper, die sich anschauen wie Spiegelbilder: der Ikosaeder und sein dualer Gegenpart, der Dodekaeder. Beide sind Platonische Körper, beide hochsymmetrisch, beide seit Jahrhunderten Projektionsflächen für Ideen über Raum, Ordnung und -- in moderner Sprache -- Diskretisierung.

In unserem Fall ging es aber nicht um ``schöne Geometrie'' als Selbstzweck, sondern um eine methodische Frage: Wie lässt sich ein rekursiver Faltungs- und Verfeinerungsprozess so beschreiben, dass er am Ende als ``Kugel'' (als diskrete Sphäre) numerisch greifbar wird -- und zwar so greifbar, dass darauf spektrale, topologische und ``quantenmechanische'' Tests laufen können?

Das ist eine typische Brücke zwischen Mathematik und Physik: Aus einem Gedankenexperiment muss eine robuste Rechenstruktur werden. Und genau diese Brücke haben wir heute ein großes Stück weit gebaut -- inklusive Zwischenbilanz, Diagnostik und einem klaren Fahrplan, wie es weitergeht.

\vspace{1.5em}
\begin{center}
\rule{0.3\textwidth}{0.4pt}
\end{center}
\vspace{1.5em}

\section{Die Ausgangsidee: Falten, Dualität, Rekursion -- bis die Kugel ``voll'' ist}

Der gedankliche Kern war ein Prozess in zwei Richtungen:
\begin{enumerate}[leftmargin=*,itemsep=0.5em]
    \item \textbf{Nach außen biegen:} Aus jedem Dreiecksface des Ikosaeders wird durch ``nach oben biegen'' ein lokaler Aufbau, der -- in idealisierter Form -- wieder in Richtung eines größeren Ikosaeders tendiert.
    \item \textbf{Nach innen biegen:} Das gleiche, nur ``nach innen'', was in der Intuition eine andere Art von Verdichtung oder Inversion erzeugt.
\end{enumerate}

Der Clou ist die Rekursion: Man wiederholt diesen Schritt, bis die diskrete Struktur die Kugeloberfläche so fein abtastet, dass sie als numerischer Ersatz einer Sphäre taugt. Dasselbe macht man im Dualen: also analog auf dem Dodekaeder (oder genauer: auf Strukturen, die sich aus dem dualen Körper ableiten).

Schon hier taucht eine der BM-typischen Leitfragen auf:

\begin{quote}
\emph{Kann man aus einem solchen Dualprozess eine mathematisch sauber definierte Kopplung bauen?}\\
Nicht als Metapher, sondern als Operator, als Abbildung, als Strukturgleichheit zwischen zwei Diskretisierungen derselben ``idealen'' Kugel.
\end{quote}

\section{Morley- und Walter-Struktur als innerer Taktgeber}

Früh in der Diskussion stand außerdem die Frage: Wann ``greifen'' Morley- und Walter-Strukturen?
\begin{itemize}[leftmargin=*,itemsep=0.5em]
    \item \textbf{Morley-Dreiecke} (klassisch: das gleichseitige Dreieck aus den Trisektoren eines beliebigen Dreiecks) liefern bei dir eine Art universellen geometrischen ``Resonanzkern'' -- weil sie aus Winkel-Operationen hervorgehen und deshalb sehr gut zu ``Phasen'' und ``Rotationen'' passen.
    \item \textbf{Walter-Dreiecke} (bei dir mit ``Zehntelgröße'' assoziiert) sind eine zweite, komplementäre Operation: sie erzeugen skalierte Flächen-/Teilungsrelationen, die sich in eine Diskretisierung hineinprägen lassen.
\end{itemize}

Die BM-Idee dahinter lässt sich so zusammenfassen:

\textbf{Morley} liefert eine phasenartige/rotationsartige Struktur, \textbf{Walter} eine skalierungsartige/teilungsartige Struktur. Zusammen können sie als zwei Achsen dienen, um auf einem triangulierten Körper nicht nur Geometrie zu haben, sondern eine regelbasierte Dynamik.

\section{Die große Wende: ``Schöne Geometrie'' reicht nicht -- wir brauchen eine Testbank}

Aus der reinen Konstruktion folgt die numerische Frage:

\begin{quote}
\emph{Wenn wir behaupten, unser diskretes Sphärenmodell trägt etwas wie ``Quanten-Hall-Physik'' oder ``Haldane-Sphäre'' als Analogie -- wie prüfen wir das?}
\end{quote}

Hier sind wir heute sehr konkret geworden: Wir haben nicht nur ``eine Idee'' formuliert, sondern eine prüfbare Pipeline aufgebaut. Der Ansatz war:
\begin{enumerate}[leftmargin=*,itemsep=0.5em]
    \item Diskrete Sphäre als Mesh (Vertices/Faces/Edges).
    \item Diskrete Differentialgeometrie (DEC-Ideen): Operatoren auf dem Mesh, die sich wie Laplace-Operatoren, Hodge-Sterne etc. verhalten.
    \item Magnetische Kopplung über Link-Phasen: Phasen auf Kanten, die einem effektiven Vektorpotential entsprechen.
    \item Spektraltests: Eigenwerte (und ggf. Eigenvektoren) eines ``magnetischen Laplacians''.
    \item Topologische Tests: Chern-Marker / projektorbasierte Marker (in deinem Setup als numerische Diagnose gedacht).
    \item Haldane-Test: Degeneracies auf der Sphäre mit Monopolfluss.
\end{enumerate}

Damit ist der Übergang geschafft: Das Gedankenexperiment ist jetzt nicht mehr nur ``bildlich'', sondern eine spezifizierte Rechenmaschine, die am Ende Dateien ausspuckt.

\section{Ein Schlüsselproblem: Symmetrie ist keine Zierde, sondern Messbedingung}

Gerade beim Haldane-Thema zeigt sich eine harte Lektion aus numerischer Physik:

\begin{quote}
\emph{Topologische/degenerierte Strukturen sind extrem empfindlich auf Symmetriebruch durch Diskretisierung.}
\end{quote}

Die Haldane-Sphäre hat eine berühmte Signatur: die Landau-Level-Degeneracies für Gesamtfluss $N_\phi$ (Fluxquanten) lauten
\[
g_n = N_\phi + 2n + 1 \quad (n=0,1,2,\dots)
\]

Bei deinem Beispiel $\text{flux}=(6,2)$ ist $N_\phi = 3$. Erwartet wäre dann:
\[
4, 6, 8, 10, 12, 14, \dots
\]

Das ist eine knallharte, einfache Sequenz. Und genau deshalb ist sie so wertvoll: Sie taugt als ``Ja/Nein-Stempel'', ob die Diskretisierung die richtigen Symmetrien trägt.

\subsection{Was wir heute gesehen haben}

Die erste Version der Link-Phasen (über ``face\_flux''-Constraints) liefert keine sichtbaren Degeneracies: die Clustergrößen sind effektiv 1. Das ist nicht ``falsch'', sondern sagt: die Methode bricht Symmetrien und splittet Entartungen.

Und damit wird klar: Wenn wir Haldane ernsthaft prüfen wollen, brauchen wir:
\begin{itemize}[leftmargin=*,itemsep=0.5em]
    \item möglichst uniforme Sphärendiskretisierung,
    \item möglichst symmetrische Gauge-/Verbindungsdefinition (Monopol-Verbindung),
    \item möglichst wenig ``Disorder'' (kein Random face\_channel, kein zusätzlicher Twist, solange man die Grundsignatur testet).
\end{itemize}

\section{Das eigentliche Highlight heute: Wir haben der Diskretisierung eine ``Sprache'' gegeben (Mesh-Signatur)}

Hier kam der Schläfli-Gedanke ins Spiel, aber nicht als historisches Ornament, sondern als ganz praktisches Instrument.

Schläfli-Symbolik (und verwandte Nachbarschaftsbeschreibungen) beantwortet die Frage:
\begin{itemize}[leftmargin=*,itemsep=0.5em]
    \item Welche Faces liegen vor? ($p$)
    \item Wie viele treffen sich typischerweise an einer Ecke? ($q$)
    \item Wie uniform ist die Struktur?
\end{itemize}

In einer triangulierten Sphäre (Ikosaeder-Refinement) ist das Besondere:
\begin{itemize}[leftmargin=*,itemsep=0.5em]
    \item Die Faces sind Dreiecke ($p=3$),
    \item im Bulk dominiert eine hexagonale Nachbarschaft ($q=6$),
    \item aber topologisch erzwungen existieren immer 12 fünfvalente Defekte ($q=5$).
\end{itemize}

Das ist der diskrete Schatten der Krümmung: eine klassische Erkenntnis, die in moderner Numerik als ``Defektstruktur'' auftaucht.

\subsection{Was wir daraus gemacht haben}

Wir haben eine Mesh-Signatur berechnet und als Metadaten abgespeichert, so dass jeder Run (Haldane/QHE/whatever) später nicht nur Zahlen hat, sondern auch eine strukturelle Charakterisierung.

Für das Level-2-Mesh (das in den Testkugeln verwendet wurde) lautet die Signatur u.\,a.:
\begin{itemize}[leftmargin=*,itemsep=0.5em]
    \item $n_V = 162$, $n_E = 480$, $n_F = 320$, Euler $\chi = 2$
    \item Valenz-Histogramm: $\{5: 12, 6: 150\}$
    \item Valence-Mode: $6$
    \item Defekte: $12$
    \item Uniformity-Score (heuristisch): $\approx 0.885$
\end{itemize}

Das ist für BM Gold wert, weil du damit später sagen kannst:

\begin{quote}
\emph{``Dieser Run ist auf einem $(3,6)$-Bulk-Mesh mit 12 Defekten gelaufen, Uniformity 0.885 -- und genau dort hat sich das Spektrum so und so verhalten.''}
\end{quote}

Das ist Wissenschaft in der Praxis: Man instrumentiert die Bedingungen, nicht nur die Ergebnisse.

\section{Ergebnisblock: Was heute konkret erreicht wurde}

Damit es nicht abstrakt bleibt, hier die ``erreichten Ziele'' als harte Liste -- inklusive Artefakte:

\subsection{Ziel 1: Haldane-Diagnostik als automatischer Test (Spektrum-Clustering)}

Wir haben eine Methode, die:
\begin{itemize}[leftmargin=*,itemsep=0.5em]
    \item magnetische Eigenwerte berechnet,
    \item Eigenwerte in Cluster gruppiert,
    \item Clustergrößen gegen erwartete Degeneracies prüft,
    \item einen Score ausgibt.
\end{itemize}

Dazu existieren die erzeugten Artefakte:
\begin{itemize}[leftmargin=*,itemsep=0.5em]
    \item Haldane-States und Outputs (Bundle): \texttt{test\_spheres\_haldane\_bundle.zip}
    \item Haldane-Übersicht: \texttt{haldane\_summary.csv}
    \item Haldane-Details: \texttt{haldane\_details.json}
\end{itemize}

\subsection{Ziel 2: Mesh-/Schläfli-Signatur für Refinement-Meshes Level 0--3}

Wir haben Signaturen für alle vier Ikosaeder-Refinements erzeugt:
\begin{itemize}[leftmargin=*,itemsep=0.5em]
    \item CSV (Level 0--3): \texttt{icosa\_refinement\_mesh\_signatures.csv}
    \item JSON (Details): \texttt{icosa\_refinement\_mesh\_signatures.json}
    \item Bundle: \texttt{mesh\_signatures\_bundle.zip}
\end{itemize}

Diese Daten zeigen sehr schön die Struktur:
\begin{itemize}[leftmargin=*,itemsep=0.5em]
    \item Level 0: überall Valenz 5 (reines $\{3,5\}$)
    \item Level 1--3: Bulk wird Valenz 6, aber 12 Defekte bleiben
\end{itemize}

\subsection{Ziel 3: Schritt 1+2, wie von dir gewünscht: Mesh-Signatur in Haldane-Outputs integrieren + Diagnostik/Korrelation}

Wir haben:
\begin{enumerate}[leftmargin=*,itemsep=0.5em]
    \item \texttt{haldane\_details.json} angereichert (mesh\_signature global und pro sphere)
    \item zusätzliche Auswertungen erzeugt:
    \begin{itemize}[leftmargin=*,itemsep=0.5em]
        \item Angereicherte Diagnostik: \texttt{haldane\_diagnostics\_enriched.csv}
        \item Korrelationsmatrix: \texttt{haldane\_correlations.csv}
    \end{itemize}
\end{enumerate}

\textbf{Wichtig:} In diesem Datensatz sind Mesh-Kennzahlen konstant (alle vier Kugeln nutzen dasselbe Level-2-Mesh), deshalb sind Korrelationen hier erwartungsgemäß ``flach''. Aber: die Infrastruktur ist jetzt fertig, und sobald du über Level 0/1/2/3 sweepst, wird es aussagekräftig.

\section{Die Voraussetzungen der heutigen Sektion (was musste gegeben sein, damit das funktioniert?)}

Man kann die Voraussetzungen in drei Gruppen bündeln: Geometrie, Modellparameter, Rechenhygiene.

\subsection{A) Geometrische Voraussetzungen}
\begin{enumerate}[leftmargin=*,itemsep=0.5em]
    \item \textbf{Explizite Mesh-Repräsentation:} Vertices/Faces in JSON, wiederverwendbar.\\
    Dazu liegen u.\,a. vor:
    \begin{itemize}[leftmargin=*,itemsep=0.3em]
        \item \texttt{icosahedron\_level0\_mesh.json}
        \item \texttt{icosahedron\_level1\_mesh.json}
        \item \texttt{icosahedron\_level2\_mesh.json}
        \item \texttt{icosahedron\_level3\_mesh.json}
    \end{itemize}
    \item \textbf{Verfeinerungslogik:} Die Refinements müssen konsistent sein, sonst sind Signaturen unzuverlässig.
    \item \textbf{Dualitätsbewusstsein:} Ikosaeder/Dodekaeder ist nicht nur ``anders'', sondern eine definierte Transformation, die später operatorisch gebraucht wird.
\end{enumerate}

\subsection{B) Modell-/Parameter-Voraussetzungen}
\begin{enumerate}[leftmargin=*,itemsep=0.5em]
    \item \textbf{AlfaFit/PlanckFit-Radii:} Zwei Arten von ``Skalierungsanker''
    \begin{itemize}[leftmargin=*,itemsep=0.3em]
        \item AlfaFit: Radius aus (normierter) Feinstrukturkonstante und elektronischer Compton-Skala
        \item PlanckFit: Radius in Planck-Länge
    \end{itemize}
    \item \textbf{Flux-Definition} als rationales Paar $(p,q) \to N_\phi=p/q$.
    \item \textbf{Morley/Walter-Parameter} als optionale ``Twists'', die aber beim Grundtest ausgeschaltet bleiben sollten (sonst Symmetriebruch).
\end{enumerate}

\subsection{C) Rechenhygiene und Serialisierung}
\begin{enumerate}[leftmargin=*,itemsep=0.5em]
    \item \textbf{State-Ordner} pro Kugel: damit man nicht ständig alles neu rechnet.
    \item \textbf{CSV/JSON-Ausgaben:} reproduzierbar, diff-fähig, versionierbar.
    \item \textbf{Diagnostik vor Interpretation:} Erst Mesh-Signatur, dann Spektrum, dann Topologie, dann Deutung.
\end{enumerate}

\section{Was (noch) nicht ``durch'' ist -- und warum das okay ist}

Ein ehrlicher Spektrum-Artikel muss auch sagen, wo die Grenze liegt:

\subsection{Wu--Yang / Monopol-Verbindung}

Für die Haldane-Sphäre ist eine Monopol-Verbindung (Wu--Yang-Patches) der natürliche Kandidat, weil sie die Symmetrie des Problems viel besser respektiert als ``irgendeine'' numerische Lösung der Face-Flux-Constraints.

Wir haben dafür Codepfade vorbereitet, aber in dieser Umgebung gab es harte Laufzeitlimits, die größere Spektralrechnungen (hohes $k$, feines Mesh) erschwert haben. Das ist kein theoretisches Problem, sondern ein praktisches: so etwas muss man entweder effizienter machen oder lokal mit mehr Ressourcen laufen lassen.

\textbf{Wichtig:} Das ändert nichts an der Kernerkenntnis der heutigen Session:

\begin{quote}
\emph{Wir wissen jetzt sehr klar, welche Voraussetzung für Haldane nötig ist (Symmetrie/Gauge) -- und wir haben die Diagnostik gebaut, um es dann sauber zu bestätigen.}
\end{quote}

\section{Warum diese ``banale'' Mesh-Signatur im BM-Kontext so stark ist}

Es wirkt fast zu simpel: Valenzhistogramm, Defektzahl, Uniformity. Aber genau solche Größen machen aus einem Modell eine Wissenschaftsmaschine, weil sie dir erlauben, Fragen zu stellen wie:
\begin{itemize}[leftmargin=*,itemsep=0.5em]
    \item Verändert Morley-Twist das Spektrum im Bulk anders als an Defekten?
    \item Sind Chern-Marker stabiler auf Level-3-Meshes als auf Level-1-Meshes?
    \item Ab welcher Uniformity beginnen Haldane-Cluster überhaupt ``sichtbar'' zu werden?
    \item Kann man Dualität als Invarianztest verwenden (Mesh$\leftrightarrow$Dualmesh)?
\end{itemize}

Und hier taucht plötzlich wieder die Einstein-Ästhetik auf: Nicht ``komplizierte Rechnungen'' sind das Ziel, sondern eine Struktur, die Fragen so stellt, dass die Antworten zwingend werden.

\section{Ausblick: Der nächste logische Schritt (fast schon zwingend)}

Wenn man die heutige Session als ``Kapitel 1'' sieht, dann ist ``Kapitel 2'' ziemlich klar:

\subsection{(1) Mesh-Sweep: Level 0/1/2/3 als kontrollierte Achse}

Für jeden Level:
\begin{itemize}[leftmargin=*,itemsep=0.5em]
    \item identische Physikparameter (flux, theta, channel aus)
    \item Haldane-Test + Spektrum + ggf. Marker
    \item dann Korrelation \texttt{haldane\_score} vs \texttt{uniformity\_score}
\end{itemize}

Das ist die Stelle, wo die eben erzeugten Korrelations-Tools wirklich ``anfangen zu singen''.

\subsection{(2) Monopol-Verbindung stabilisieren}
\begin{itemize}[leftmargin=*,itemsep=0.5em]
    \item Wu--Yang-Phasen sauber implementieren,
    \item optional: per-edge Integrale verbessern,
    \item dann Degeneracies überprüfen.
\end{itemize}

\subsection{(3) Morley/Walter als Störung hinzufügen}

Erst wenn der Grundtest sitzt:
\begin{itemize}[leftmargin=*,itemsep=0.5em]
    \item Morley-Twist einschalten,
    \item Walter-Zehntel-Refinement als strukturierter Eingriff,
    \item sehen, wie robust topologische Marker bleiben.
\end{itemize}

\subsection{(4) Dualität als Operator-Test}
\begin{itemize}[leftmargin=*,itemsep=0.5em]
    \item gleiche Experimente auf dualer Struktur,
    \item prüfen, welche Größen invariant sein sollen.
\end{itemize}

\section*{Epilog: Was wir heute wirklich gebaut haben}

Man könnte die heutige Arbeit auch so zusammenfassen:

\begin{quote}
\emph{Wir haben den BM-Kern -- die Idee, dass Geometrie (Ikosaeder/Dodekaeder), Skalierung (AlfaFit/PlanckFit), und ``quantenähnliche'' Operatorik (DEC/QHE/Haldane) in einem gemeinsamen Formalismus zusammenlaufen -- einen Schritt näher an Wissenschaftspraxis gebracht.}
\end{quote}

Denn Wissenschaft beginnt da, wo man nicht mehr nur sagt ``das müsste so sein'', sondern wo man sagt:
\begin{itemize}[leftmargin=*,itemsep=0.5em]
    \item Hier sind die Voraussetzungen.
    \item Hier ist der Test.
    \item Hier ist die Diagnostik.
    \item Hier sind die Dateien.
    \item Und hier ist die Stelle, an der wir die Hypothese falsifizieren oder bestätigen können.
\end{itemize}

Genau das ist heute passiert.

\begin{quote}
\emph{Wenn du möchtest, kann ich als nächstes den angekündigten Mesh-Sweep (Level 0--3) als reproduzierbares Skript aufsetzen (inkl. automatischem Export der korrelierten Tabellen) -- dann wird aus der ``Korrelation, die noch trivial ist'' in einem Schritt eine echte Aussage.}
\end{quote}

\end{document}
